% Generated draft: requires mathematical and source review.
\documentclass[11pt]{article}
\usepackage[margin=25mm]{geometry}
\usepackage{fontspec}
\setmainfont{Latin Modern Roman}
\usepackage{amsmath,amssymb,mathtools,unicode-math}
\setmathfont{Latin Modern Math}
\usepackage{enumitem,needspace,xurl,hyperref}
\setlength{\parindent}{0pt}
\setlength{\parskip}{5pt}
\setlength{\emergencystretch}{3em}
\newcommand{\N}{\mathbb{N}}
\newcommand{\Z}{\mathbb{Z}}
\newcommand{\Q}{\mathbb{Q}}
\newcommand{\R}{\mathbb{R}}
\newcommand{\C}{\mathbb{C}}
\newcommand{\F}{\mathbb{F}}
\newcommand{\A}{\mathbb{A}}
\newcommand{\PP}{\mathbb{P}}
\newcommand{\E}{\mathbb{E}}
\newcommand{\eps}{\varepsilon}
\newcommand{\dd}{\,\mathrm{d}}
\begin{document}
\section*{UnsolvedMath ID: 20002450}
\textbf{Source problem code:} catalogue entry
\subsection{Definitions and mathematical statement}
\textbf{Source question with missing antecedent.} Can one construct an Evans function for ``this problem'', and how does it behave when the absolute spectrum touches the essential spectrum and the usual exponential dichotomy is unavailable? The supplied extraction contains no underlying operator, wave, spectral ODE, asymptotic matrices, function space or definition of the relevant spectra. The following heading, P. Howard, ``Combination structures in viscous conservation laws'', is a boundary to another contribution, not a recovered attribution or operator for this question. Spectral contact is retained as a source assertion at the relevant threshold, not as a general claim that no dichotomy exists at any parameter.

The aid refers to an omitted ``stationary \(p=5\) gKdV'' setup and reports an inconsistency: its first integral excludes a nonzero smooth zero-background homoclinic. It also reports that a standard co-moving repair has exponentially decaying \(L^2\) eigenfunctions, including at zero or embedded spectral values, and so cannot supply the merely algebraic mode intended there. Neither the printed equation, repaired operator nor far-field cubic is supplied in this record. Those claims cannot be promoted to precise gKdV theorems or used to choose the missing antecedent silently.

\textbf{Reported rank-one threshold model.} A separate explicitly conditional model is the scalar channel, on an outward half-line coordinate \(s\geq0\),
\[
 z'(s)=\left[-\mu(\lambda)-a(\lambda)(1+s)^{-\beta}
                   +r(s,\lambda)\right]z(s),\qquad \beta>0.
\]
Assume \(\mu,a\) and the remainder are analytic in the spectral parameter on a stated domain, with the \(L^1\) bound and tail convergence of \(r(\cdot,\lambda)\) locally uniform there. These uniformity conditions are necessary for the analytic limiting assertion. Put
\[
 H_\beta(R)=\int_0^R(1+s)^{-\beta}\,\mathrm ds
 =\begin{cases}((1+R)^{1-\beta}-1)/(1-\beta),&\beta\ne1,\\
 \log(1+R),&\beta=1.\end{cases}
\]
For a nonzero mode,
\[
 z(s)=z(0)\exp\left[-\mu s-aH_\beta(s)+\int_0^s r(t,\lambda)\,\mathrm dt\right].
\]
If the cutoff mode is normalized at \(R\) by \(z_R(R)=e^{-\mu R}\), then
\[
 z_R(0)=\exp\left[aH_\beta(R)-\int_0^Rr(s,\lambda)\,\mathrm ds\right].
\]
This specifies the cutoff convention; changing it changes nonzero scalar determinant factors.

For a matching determinant multilinear in independent rank-one channel columns \(\sigma\), with fixed matching vectors and all other factors combined in \(D_0(\lambda)\), the corresponding exact model factorization is
\[
 D_R(\lambda)=
 \exp\left[\sum_\sigma a_\sigma H_{\beta_\sigma}(R)\right]
 \exp\left[-\sum_\sigma\int_0^Rr_\sigma(s,\lambda)\,\mathrm ds\right]D_0(\lambda).
\]
Removing the first factor gives the analytic limit
\[
 D_{\mathrm{ren}}(\lambda)=
 \exp\left[-\sum_\sigma\int_0^\infty r_\sigma(s,\lambda)\,\mathrm ds\right]
 D_0(\lambda).
\]
The integrable remainder contributes a finite nonzero factor and must not be dropped from an exact identity. This is a conditional rank-one matching construction; it is not an Evans function for an unspecified full PDE or a proof that its zeros have a prescribed spectral multiplicity.

At a threshold \(\lambda_0\) with \(\mu(\lambda_0)=0\) and \(\beta=1\), the aid's sharp integrability test is
\[
 z\in L^2([0,\infty))\quad\Longleftrightarrow\quad
 \operatorname{Re}a(\lambda_0)>\tfrac12
\]
for a nonzero mode and an \(L^1\) remainder. A full-line mode must meet this condition on both outward channels. At equality it is not square integrable; for \(0\leq\operatorname{Re}a\leq1/2\) the bounded non-\(L^2\) channel may be called a threshold resonance under a stated convention. A growing non-\(L^2\) mode is not automatically a resonance. The model explains a possible algebraic threshold renormalization but leaves the missing original operator and its spectral interpretation unresolved.
\subsection{Short English statement}
The source does not identify the operator behind its Evans-function question. The aid's separate rank-one model exhibits an explicit cutoff factor and an algebraic threshold integrability boundary. The draft fixes the cutoff convention and retains the analytic remainder factor without claiming an Evans construction for the missing PDE.
\subsection{Sources}
\begin{itemize}
\item Local UnsolvedMath registry, permanent ID \(20002450\), catalogue entry, \texttt{aim-pdes-notes.json}, \texttt{notes[61]}, workshop \emph{Stability criteria for multi-dimensional waves and patterns}, source item 14: \url{https://aimath.org/WWN/multidimwaves/multidimwaves.pdf}.
\item Supplied extraction includes the subsequent heading P. Howard, ``Combination structures in viscous conservation laws''; it does not identify that heading as the author/setup of the preceding question.
\item Supplied research summary on a stationary/co-moving \(p=5\) gKdV audit and an independent rank-one threshold-channel renormalization. No full equations or separate authoritative citations are supplied. Attempt 1: partial\_result, status at run ``unknown'', rigor ``proved'', low novelty confidence; the registry says partially\_solved. The machine-generated aid is not peer reviewed and requires independent verification before citation. No PDF visit is claimed.
\end{itemize}
\end{document}
